\documentclass[11pt]{article}
\usepackage[margin=1in]{geometry}
\usepackage[T1]{fontenc}
\usepackage{lmodern,amsmath,booktabs,graphicx}
\usepackage[hidelinks]{hyperref}
\usepackage{xurl}
\setlength{\parindent}{0pt}
\setlength{\parskip}{6pt}
\setlength{\emergencystretch}{3em}
\graphicspath{{simulation_output/estimator_visible_trajectory_20261003/}}
\title{A Target Trajectory That Preserves LiDAR Coverage}
\author{Pan Dynamics Collision Avoidance Research}
\date{October 3, 2026}
\begin{document}
\maketitle

\section{Implemented behavior}
The target now adjusts its speed from relative longitudinal spacing instead
of receiving an independent fixed 11 m/s reference. The initial road spacing
is 18 m between actor origins. Its desired center gap oscillates smoothly
between 13 and 15 m, with zero commanded offset from its lane center.
The ego speed controller retains its previous 11 m/s setting. Parked vehicle
map decorations remain disabled.

The previous 15 m label denoted spawn spacing; its stable evaluation actually
placed the target center only 4.729--4.812 m ahead. The preceding visibility
study identified missing low vehicle surfaces at those distances. This change
addresses scenario geometry, without changing the rectangle fitter, observer,
sensor mounting, LiDAR resolution or noise. It does not claim a new estimator
RMSE result.

\section{Reference and physical controls}
Let $s$ be elapsed driving time after the 1 s stationary calibration period,
$d$ the target collision-box center projected onto the ego forward axis
relative to the ego center of mass, and $v_e$ the ego forward velocity.
The scenario policy is
\[
 r(s)=14+\sin(0.6s)\quad[\mathrm m],\qquad
 v_t^{\rm ref}=\operatorname{clip}_{[0,22]}
 \bigl(v_e+0.6\cos(0.6s)+0.8(r(s)-d)\bigr)\quad[\mathrm{m/s}].
\]
For aligned straight motion with exact, unsaturated velocity tracking,
$\dot e=-0.8e$ for $e=d-r$. Acceleration lag, saturation and turning mean this
is not a guaranteed visibility invariant. The 9--24 m configuration interval
is an evaluation envelope, not a hard state constraint.

The target follows map waypoints with the existing physical steering rule.
Its speed servo retains throttle feedforward 0.25, proportional speed gain
0.18, brake gain 0.25 and maximum brake 0.8. Maximum target throttle increases
from 0.65 to 1.0 to give the slower-launching target enough authority to
recover spacing. Throttle is zero when brake exceeds 0.05; a zero speed
reference applies full brake. Vehicles are not teleported and their velocity
is not assigned directly.

Configuration is in \path{config/carlaPerceptionScene.json}; the simulator-independent
policy and speed servo are in \path{scripts/perceptionTargetTrajectory.py}.
The local full RGB/LiDAR/GNSS/IMU capture entry point is
\path{simulation_output/estimator_visible_trajectory_20261003/captureNative.py}.
It consumes this configuration and the same policy and servo by default.
The local native coverage driver is \path{validateNative.py} in that directory.
Simulator adapters and recorded data remain outside tracked project source.
Actor state is used only for scenario control and separate scoring, never
as a perception measurement or estimator prior. The prediction-initialized,
free-pose LiDAR fitting method remains unchanged.

\section{Native validation}
One 900-frame run at 50 Hz exercised CARLA 0.10.0 on Town10HD\_Opt, spawn 105,
with a Lincoln MKZ ego and Dodge Charger target. There were 50 stationary,
150 driving-warmup and 700 evaluation frames. Evaluation $t=0$ is capture
frame 200, preserving the earlier time convention. Stable evaluation starts
at $t=2$ s, capture frame 300. The nominal recording duration is 18 s.
The route includes turns: maximum relative heading was $60.92^\circ$ and
maximum ego-frame lateral offset was 7.506 m, despite zero lane-offset command.

Native ordinary LiDAR retained 64 channels, 1.28 million points/s, 50 Hz,
80 m range, $[-30,10]^\circ$ elevation, 0.02 m noise and zero dropoff, with seed
20261003. A co-located native semantic LiDAR identifies target hits strictly
for coverage scoring. ``Low body'' here means target returns at ego-COM
$z\leq0.5$ m; it is a diagnostic count, not a new perception input.

\begin{center}\small
\begin{tabular}{@{}lrrr@{}}\toprule
Quantity & All frames & Stable & Matched old window \\
 & 900 & 600 & 200 \\\midrule
Evaluation time (s) & $-4$--13.98 & 2--13.98 & 2--5.98 \\
Center longitudinal gap (m) & 9.834--17.844 & 10.193--14.695 & 10.193--13.037 \\
Target LiDAR hits/frame & 47--202 & 72--190 & 104--190 \\
Low-body hits/frame & 16--113 & 29--109 & 44--90 \\
Frames without target hits & 0 & 0 & 0 \\
Frames without low-body hits & 0 & 0 & 0 \\
All box corners inside LiDAR & 900/900 & 600/600 & 200/200 \\
All corners inside projected image & 900/900 & 600/600 & 200/200 \\
\bottomrule\end{tabular}
\end{center}
The lowest bounding-box corner elevation over the full run was
$-16.289^\circ$, leaving $13.711^\circ$ above the lower LiDAR boundary.
All corners remained within 20.665 m range. Their front-camera geometric
projection retained at least 46.216 pixels to an image boundary. These are
true-box geometric checks, not observed YOLO detections or proof that every
vehicle surface is visible. Low-body native hits provide additional evidence
that the particular near-range truncation has been avoided.

Seven Python behavior tests passed, covering feedback direction, smooth
reference velocity, speed bounds, acceleration-lag convergence, invalid
inputs, physical actuation and configuration coverage consistency. All
850 driving throttle/brake commands from the inline native-test servo were
reproduced by the extracted default-capture helper within $2.981\times10^{-8}$,
consistent with the simulator's float precision. Sensor frame IDs matched on
all 900 ticks and consecutive frames were uninterrupted. No collision events
were recorded.

\section{Scheduling, scope and limits}
The shared port 2000 was not used. An exclusively owned server on port 2011
held both existing local coordination locks. It ran with NullRHI and no image
rendering under a 5 GiB memory limit, zero swap allowance and 250\% CPU quota.
This avoided competing for the nearly full GPU memory. The two vehicles and
three sensors received destruction requests without errors; original world
settings and 143 hidden decoration components were restored. Immediate
actor-list queries still returned counts of two vehicles and three
sensors. The owned server was subsequently stopped successfully and both its
process exit and closed ports were verified, removing its remaining resources.

Native physics and LiDAR were exercised. RGB rendering, YOLO inference and
estimator replay were not performed in this run. The full sensor capture
adapter was prepared and syntax checked; its target actuation was checked
against every executed native driving command. This finite scenario is not
a guarantee for other routes, sudden braking, occlusion, weather, sensor
mounts or arbitrarily long runs. New recordings must retain explicit coverage
checks. The summary and provenance are in
\path{report/ESTIMATOR_VISIBLE_TRAJECTORY_20261003/}.

\clearpage
\begin{figure}[p]
\centering
\includegraphics[width=\linewidth,height=0.86\textheight,keepaspectratio]{trajectory-coverage.pdf}
\caption{Observed spacing and target returns during the extended native run.
The green interval is the stable evaluation window. Physical launch lag and
turning change actual spacing, but low-body hits persist in every frame.}
\end{figure}
\end{document}
